\uuid{LafO}
\titre{Arithmétique}
\auteur{bodin}
\organisation{exo7}
\licence{CC BY-NC-SA 4.0}
\datecreate{2020-01-21}
\niveau{L1}
\module{Algèbre}
\chapitre{Arithmétique dans Z}
\sousChapitre{Autre}
\difficulte{}
\theme{Arithmétique}
\exo7qcmid{arnaud.arnaud.012}

\contenu{
\begin{question}
\qclassification{103.99}
Quelles sont les assertions vraies ?
\begin{answers}
    \good{Soit $n\in \Nn$ un entier, alors $(n+1)(n+2)(n+3)(n+4)$
est divisible par $24$.}

    \bad{Soit $n \ge 6$ un entier pair alors $\frac{n}{2}$ est impair.}

    \good{La somme et le produit de deux nombres pairs est un nombre pair.}

    \bad{$a|b$ et $a'|b' \quad \implies  \quad aa'|bb'$.}

    \bad{$a|b$ et $a'|b' \quad \implies  \quad a+a'| b+b'$.}

\end{answers}
\end{question}
}
